\chapter{Introduction and scientific background}\label{ch:introduction}

\section{Cellular responses to ionising radiation}

Ionising radiation kills cells mainly by damaging DNA, and the DNA double-strand break is
the lesion most closely tied to loss of reproductive capacity. The reference measure of how
much radiation a cell population survives is the \emph{clonogenic assay}, which counts the
cells that retain unlimited proliferative capacity and form colonies after irradiation
\parencite{puck1956}. For a dose $D$, the surviving fraction is the plating efficiency of
irradiated cells relative to that of unirradiated controls,
\begin{equation}
  \mathrm{SF}(D) \;=\; \frac{N_{\text{colonies}}(D)\,/\,N_{\text{seeded}}(D)}
                             {N_{\text{colonies}}(0)\,/\,N_{\text{seeded}}(0)} .
  \label{eq:sf}
\end{equation}
The surviving fraction at 2~Gy (SF2) is the most widely used single summary, because 2~Gy
is the conventional daily fraction in radiotherapy and because SF2 has been associated with
clinical outcome in some settings \parencite{fertil1981,west1993}.

Cell lines differ substantially and reproducibly in intrinsic radiosensitivity, and these
differences are only partly explained by tissue of origin \parencite{yard2016}. The
determinants include DNA damage response and repair capacity, cell-cycle checkpoint
control, proliferation rate, apoptotic propensity, oxygenation and redox state. Each of
these processes leaves a transcriptional footprint, which is why prediction from gene
expression is plausible in principle.

\section{The \textit{Moringa oleifera} experimental motivation}\label{sec:moringa-study}

The motivating study is an unpublished manuscript, \enquote{Dichloromethane seed and leaf
extracts of \textit{Moringa oleifera} (Lam) enhance proliferation in irradiated rodent
fibroblasts \textit{in vitro}}. The scanned copy available to the project is incomplete:
printed pages 1, 2, 4, 5, 6, 13, 14 and 15 are present, while page~3 and pages 7--12 are
missing. The plant material and extraction methods, the end of the discussion, the
conclusions and the reference list therefore cannot be read; authors, venue, date and DOI
are not shown, and no published version was located. No missing method, result or
conclusion has been reconstructed in this report.

The available pages document the following design. Crude DCM seed and leaf extracts were
reconstituted in phosphate-buffered saline, sterile-filtered and applied at
0--100~µg/mL to V79 (described as apparently normal Chinese hamster lung fibroblasts),
DU145 (human prostate cancer) and MCF-7 (human mammary adenocarcinoma). For each cell line
and extract, one plate was left unirradiated and a second received 2~Gy of X-rays at
1~Gy/min. Adherent biomass was measured 72~h later by crystal-violet optical density at
590~nm, with dose-matched untreated controls subtracted. The authors integrated the change
in optical density over extract concentration, giving an area $A_D$ at each dose $D$, and
reported a \emph{proliferation enhancement factor}
\begin{equation}
  \mathrm{PEF} \;=\; \frac{A_{2\,\mathrm{Gy}}}{A_{0\,\mathrm{Gy}}} .
  \label{eq:pef}
\end{equation}
For the seed extract, PEF exceeded 1 in all three lines (V79 1.25, DU145 1.12, MCF-7 1.16).
For the leaf extract, PEF was 1.26 in V79, 0.90 in DU145 and 0.99 in MCF-7.

Four cautions follow directly from the design, and they constrain everything that follows.
\begin{enumerate}
  \item \textbf{PEF is not SF2.} PEF is a ratio of integrated extract-versus-control
    differences in 72-hour biomass measured with the extract present; SF2 is a surviving
    fraction of untreated cells measured over one to three weeks of colony growth. The two
    are never equated, pooled or plotted on one axis in this project.
  \item \textbf{\enquote{Enhancement} denotes two different quantities in the manuscript.}
    By PEF the seed extract enhanced all three lines; by growth relative to untreated
    control only V79 increased, while the two carcinoma lines showed no change or reduced
    growth. The discussion's statement about the fibroblasts refers to the second quantity,
    the results' statement about all lines to the first.
  \item \textbf{The normal-versus-malignant comparison is confounded.} The only apparently
    normal line is also the only non-human line, so species, tissue of origin, phenotype
    and malignancy vary together. Only the leaf extract shows a normal/malignant difference
    in PEF.
  \item \textbf{The design cannot support a selectivity claim.} There is a single non-zero
    dose and no dose--response; the replicates are technical wells rather than independent
    biological repeats; dose is confounded with plate; and no vehicle or cell-free
    interference control is described. The available pages contain no clonogenic,
    cytotoxicity, apoptosis, DNA-damage or migration assay, no normal human cell line and
    no mechanistic experiment.
\end{enumerate}
The \textit{Moringa} measurements are therefore used as \emph{biological motivation only}.
They are never used as training, calibration or validation data, and no claim of selective
phytochemical radioprotection is made anywhere in this report.

\section{Why cellular radiosensitivity is hard to predict}

Published attempts to predict cell-line radiosensitivity from expression have performed
poorly under honest evaluation. The most directly relevant prior work, RadSigBench,
re-fitted seven published radiosensitivity signatures and benchmarked them against
resampled random gene sets, reporting a mean out-of-sample $R^2$ of 0.01 on CCLE data and
0.00 on NCI-60, with the published signatures performing no better than random gene sets
\parencite{oconnor2022}. This is the single most important prior for the present work: it
sets the expectation that the signal is weak, so a near-zero $R^2$ here is a replication
rather than a surprise.

Three difficulties recur. The first is \emph{sample size}: panels with reliable clonogenic
data number in the hundreds, against tens of thousands of candidate molecular features. The
second is \emph{label noise}: literature-aggregated SF2 values combine laboratories, dose
rates and protocols, and an audit of 1,672 radiobiology papers found that 30.5\% did not
report replicates and 32.3\% did not report dose rate \parencite{oike2020}. The third is
\emph{leakage through relatedness}: isogenic derivatives, sublines and cross-contaminated
lines share molecular profiles, so a split that separates records rather than biological
entities inflates apparent performance.

\section{The case for public radiogenomic data}

A three-line extract experiment cannot address intrinsic radiosensitivity across cell
types, but public resources can. CLARIFi aggregates literature clonogenic radiosensitivity
\parencite{clarifi2024}; DepMap provides harmonised baseline expression for the same lines
\parencite{depmap24q4}; Cellosaurus provides authoritative identity, problem flags and
relatedness \parencite{bairoch2018}; and the MSigDB Hallmark collection provides 50 compact,
interpretable pathway summaries \parencite{liberzon2015}. Together they allow a properly
grouped, prespecified benchmark on a defensible endpoint.

\section{The value and the limits of interpretable machine learning}

Interpretable models matter here for two reasons: a pathway-level summary is biologically
meaningful, and feature attributions can be examined for stability. Attribution, however,
has hard limits. Permutation importance describes a \emph{fitted model}, not biology;
correlated features share or exchange importance; and the size of an importance is not a
molecular effect size. Attribution is only as trustworthy as the prediction it explains,
which is why every pathway statement in this report is capped at an exploratory tier while
the prediction gain remains small and unvalidated.

\section{Rationale for prioritising candidate compounds}

The original study raises a question that the machine-learning benchmark alone cannot
answer: if cells differ in radiation response, and some of that difference is reflected in
pathway activity, which \textit{Moringa} constituents are worth testing? The answer has to be
built compound by compound from independently published chemistry and biology, with the
extraction solvent tracked explicitly, because the extracts used in the original study were
never chemically characterised.
